\section{Agentic Construction}
\label{sec:agentic}

\subsection{What ``agentic'' means here, and what it does not}

The construction ran in a coding-agent runtime: a language model with filesystem,
shell and database access, directed by one human owner, writing and executing code as
well as reading sources. It did \emph{not} use ReAct, AutoGen, LangGraph, CrewAI,
Reflexion or any comparable orchestration library; we searched the project's
dependency manifests and found none, and we mention the families only because the
related-work section discusses them (\cref{sec:related}). Orchestration here is the
project's own artefact protocol --- briefs, staging directories, proposal files,
gates, rulings and receipts --- rather than a framework.

Two consequences follow. The workflow is inspectable, because each step left a file.
And it is not a reusable system: what generalises is the protocol, not an
implementation.

\begin{figure}[tbp]
\centering
\resizebox{0.86\textwidth}{!}{% Figure: the construction pipeline, with the actor for each stage in the left
% gutter and what the stage may not do in the right gutter. Reconstructed from the
% repository's own artefact protocol; every stage names a file family that exists.
\begin{tikzpicture}[
  font=\footnotesize,
  every node/.style={align=center},
  st/.style={draw=black!50, rounded corners=2.5pt, inner sep=4.5pt, fill=white,
             text width=44mm, minimum height=7mm, line width=0.55pt},
  det/.style={st, fill=SeaGreen!10, draw=SeaGreen!60!black},
  agn/.style={st, fill=SkyBlue!14, draw=SkyBlue!65!black},
  adv/.style={st, fill=BurntOrange!11, draw=BurntOrange!70!black},
  own/.style={st, fill=RoyalPurple!9, draw=RoyalPurple!60!black},
  fin/.style={st, fill=black!6, draw=black!60, line width=0.8pt},
  fl/.style={-{Latex[length=2.1mm,width=1.6mm]}, draw=black!60, line width=0.7pt},
  who/.style={font=\scriptsize\bfseries, text=black!58, anchor=east,
              text width=26mm, align=right},
  nope/.style={font=\scriptsize, text=black!58, anchor=west, text width=52mm,
               align=left},
  hd/.style={font=\scriptsize\bfseries, text=black!52},
]

\node[hd] at (-5.1,1.15) {actor};
\node[hd] at (0,1.15) {stage};
\node[hd] at (5.3,1.15) {what the stage may not do};

\node[det] (s1) at (0,0.25)   {\textbf{Source acquisition}\\[-1pt]
                               \tiny hashed snapshot; rights recorded per file};
\node[det] (s2) at (0,-1.45)  {\textbf{Canonical corpus + identity}\\[-1pt]
                               \tiny \mbox{UUIDv5} over a \mbox{URN};
                               byte-reproducible build};
\node[agn] (s3) at (0,-3.55)  {\textbf{Parallel layer proposal}\\[-1pt]
                               \tiny one staging directory per domain\\[2pt]
                               \tiny\itshape attribution \textperiodcentered\ textual
                               reuse \textperiodcentered\ formulae
                               \textperiodcentered\ entities
                               \textperiodcentered\ ritual\\[-1pt]
                               \tiny\itshape translation \textperiodcentered\ notation
                               \textperiodcentered\ audio\\[2pt]
                               \tiny eight layers, proposed independently;\\[-1pt]
                               \tiny each lands in staging, none in the graph};
\node[det] (s4) at (0,-5.35)  {\textbf{Staging validation}\\[-1pt]
                               \tiny the contract reports its own coverage};
\node[adv] (s5) at (0,-7.05)  {\textbf{Gates A / B / C}\\[-1pt]
                               \tiny integrity, conformance, attack};
\node[det] (s6) at (0,-8.75)  {\textbf{Grading}\\[-1pt]
                               \tiny layer, tier, precision, surface ---
                               applied graph-wide, last};
\node[own] (s7) at (0,-10.55) {\textbf{Plan $\to$ dry run $\to$ \textsc{go}
                               $\to$ import $\to$ readback}};
\node[det] (s8) at (0,-12.45) {\textbf{Invariants and scorecard}\\[-1pt]
                               \tiny 20 invariants, 11 gates, all at zero};
\node[fin] (s9) at (0,-14.15) {\textbf{VedGraph}\\[-1pt]
                               \tiny 165{,}737 nodes, 510{,}905 relationships};
\node[fin] (s10) at (0,-15.65){\textbf{VedAnvaya}\\[-1pt]
                               \tiny read-only \mbox{API} and interface};

\foreach \a/\b in {s1/s2, s2/s3, s3/s4, s4/s5, s5/s6, s6/s7, s7/s8, s8/s9, s9/s10}
  \draw[fl] (\a) -- (\b);


% ---- left gutter: actor ----
\node[who] at (-2.45,0.25)    {deterministic};
\node[who] at (-2.45,-1.45)   {deterministic};
\node[who] at (-2.45,-3.35)   {\textbf{agent}\\proposes};
\node[who] at (-2.45,-5.35)   {deterministic};
\node[who] at (-2.45,-7.05)   {\textbf{a second agent}\\attacks};
\node[who] at (-2.45,-8.75)   {deterministic};
\node[who] at (-2.45,-10.55)  {\textbf{owner} \textsc{go};\\code verifies};
\node[who] at (-2.45,-12.45)  {deterministic};
\node[who] at (-2.45,-14.15)  {---};
\node[who] at (-2.45,-15.65)  {---};

% ---- right gutter: the refusal at each stage ----
\node[nope] at (2.45,0.25)
  {a permissive licence is never inferred, only evidenced; absence of a licence
   field is not a grant};
\node[nope] at (2.45,-1.45)
  {no source, text, script, edition, build or sequence position may enter an
   identifier};
\node[nope] at (2.45,-3.35)
  {a model may write exactly one status, \code{CANDIDATE}; the predicate
   vocabulary is closed and its refusals are named};
\node[nope] at (2.45,-5.35)
  {a check that skipped its population must show as a coverage shortfall, not as
   a pass};
\node[nope] at (2.45,-7.05)
  {Gate~C is a separate file from Gate~B, ``so the same codepath cannot validate
   itself'', and carries a negative control};
\node[nope] at (2.45,-8.75)
  {an unlisted relationship type \emph{raises}; the fall-through default is the
   weakest tier, never the strongest};
\node[nope] at (2.45,-10.55)
  {nothing is imported until a readback; an importer's own success report is not
   evidence};
\node[nope] at (2.45,-12.45)
  {the data may not declare itself valid --- a gate whose declared set includes
   the graph cannot fail};
\node[nope] at (2.45,-14.15)
  {every edge carries its epistemic layer, so a consumer can exclude
   model-derived claims with a predicate};
\node[nope] at (2.45,-15.65)
  {typing is rendered where the reader is, not on a methods page};

\end{tikzpicture}
}
\caption{The construction pipeline. The left gutter names the actor for each stage;
the right gutter names what that stage is forbidden to do, which is where the
method actually lives. Agents propose into staging; deterministic code validates,
grades and verifies; the owner issues the one decision that lets a mutation run.}
\label{fig:pipeline}
\end{figure}

\subsection{The construction pipeline}

\Cref{fig:pipeline} shows the pipeline the artefacts describe. Its spine is the same
for every layer: a source is snapshotted and hashed; a canonical corpus is built
deterministically and given identity; a layer is proposed into a \emph{staging}
directory that is validated and gated before any graph write; the owner issues a
\textsc{go} against a hashed plan; the import executes; and a readback compares what
was sent with what landed.

The staging step is the load-bearing one. A layer's proposal is a directory of
row-level JSONL plus a manifest, validated by a contract that reports \emph{its own
coverage} --- the validator states how many rows each check evaluated, so a check
that silently skipped its population is visible as a coverage shortfall rather than
as a pass. During the completeness campaign the rule was stated flatly: nothing is
imported until a readback.

\subsection{The import protocol}
\label{sec:import}

Every graph mutation in the later campaigns followed a five-step protocol, and its
steps exist because each was once skipped:

\begin{enumerate}[leftmargin=1.9em,itemsep=1pt,topsep=2pt]
\item \textbf{Plan.} A hashed artefact enumerating the exact nodes, edges and
property writes, with the expected census delta.
\item \textbf{Dry run.} The plan re-computed against the live graph, emitting a
\textsc{go} or the reason it cannot. Dry runs were regenerated whenever the tree
changed, because a dry run computed against a different state is not a dry run.
\item \textbf{Owner \textsc{go}.} A recorded human decision against a specific plan
hash.
\item \textbf{Execute.}
\item \textbf{Readback.} An independent re-measurement comparing rows \emph{sent}
against rows \emph{landed}, with the census delta checked against the promise.
\end{enumerate}

The protocol's value is measurable in what it caught: a first import that
over-wrote by 5{,}930 edges, a second that left 12{,}033 nodes unreachable, a third
that landed 13{,}187 elements the domain rules excluded, and one that silently
overwrote 102 curated display labels. Roughly nine plan/dry-run/\textsc{go} cycles
preceded one clean readback in the largest wave.

It also failed in an instructive way. One readback printed \code{READBACK\_CLEAN}
having compared zero rows against zero rows, and a separate import's own readback
reported success while two graph-wide contracts were unmet --- the scorecard, not the
importer, found 1{,}136 edges that were ungraded, unlabelled and unmarked. The
lesson the repository records is the one we would generalise: \emph{an importer's own
success report is not evidence}.

\subsection{The multi-agent structure, as the artefacts show it}
\label{sec:multiagent}

The clearest instance is the final closure sprint, and we describe it using the
repository's own names rather than inventing roles.

Forty-three remaining gaps were partitioned across \textbf{seven domain agents}, one
owner each, under a stated rule: every implementation-fixable row has exactly one
owner agent, and no row is unassigned. The agents are named by domain ---
\code{AGENT\_1\_FORMULA\_CROSS\_VEDA} (3 gaps),
\code{AGENT\_2\_ENTITY\_DEVATA\_IDENTITY} (12),
\code{AGENT\_3\_MORPHOLOGY\_SEMANTICS} (7), \code{AGENT\_4\_RITUAL} (4),
\code{AGENT\_5\_QUALITY\_REFERENCE} (6), \code{AGENT\_6\_PRODUCT\_SURFACE} (5) and
\code{AGENT\_7\_TRANSLATION\_SAMAVEDA} (6). Each produced measurements, row-level
evidence and a file named \code{proposed\_mutations.json}. \emph{None of them wrote
to the graph.}

An eighth agent had no gaps assigned. Its artefacts are an audit of the other seven,
its recorded database access is \code{READ\_ONLY}, and it is the adversary in the
pattern. A \textbf{lead} then reconciled proposals and findings into an integration
plan, a set of rulings and a receipt.

\Cref{fig:adversarial} shows the resulting cycle. Its most important property is
visible in one ruling. Agent~3 proposed re-anchoring a relationship so that role
fillers hang from semantic assertions. Agent~8 measured that the change would create
2{,}052 endpoint-signature violations, because the declared signature expects a
passage endpoint and zero of the assertion nodes carry that label. The lead verified
the measurement independently, and then ruled neither ``approved'' nor ``rejected''
but \emph{approved with the ontology moving in the same change} --- because the
migration card had declared the new signature all along, and either half alone would
leave 2{,}052 violations. A test now pins the signature to the card.

Three features of that ruling recur across the corpus of rulings and are, we think,
the transferable part:

\begin{itemize}[leftmargin=1.9em,itemsep=1pt,topsep=2pt]
\item \textbf{The adversary measures rather than opines.} Agent~8's findings are
counts against the live store, so they can be re-run. One of them --- that a
registry-status check ``only compared a graph measurement to a pre-declared integer;
it never read the status field'' --- is a defect in a \emph{check}, not in data.
\item \textbf{The integrator re-measures before ruling.} The lead's rulings record
\code{lead\_measurement} separately from the agent's claim, and in one case correct
the agent's figure: an agent measured the live layer at 4{,}865 unreviewed
assertions while its own staged delta would take the post-integration value to
35{,}139 --- a prediction, against the 35{,}131 the layer holds today.
\item \textbf{Trades are recorded, not netted.} One ruling notes that closing a
coverage gap multiplies another gap's unreviewed population sevenfold and deepens a
third, and records the trade instead of reporting the improvement alone.
\end{itemize}

The adversary also found defects nobody had a gap for. It measured that the exported
public artefact under-declared its own inputs --- 31 of 36 exported labels
undeclared, 129{,}691 edges and 32{,}044 node rows invisible to the hash that was
supposed to cover them, and 4{,}016 public nodes dropped silently. And it found a
user-visible search defect by querying rather than by reading code: 5{,}123 of 5{,}839
search derivatives in one corpus carried a private-use sentinel character, so a
reader searching for \iast{\d{r}tasya} got zero hits while the sentinel form got 94.

\begin{figure}[tbp]
\centering
\resizebox{\textwidth}{!}{% Figure: propose -> attack -> adjudicate, as the final closure sprint's artefacts
% record it. Names and counts are the repository's own (gap_assignment.json,
% agent8/audit_findings.json, lead/rulings.json); no role is invented to make the
% diagram symmetrical.
\begin{tikzpicture}[
  font=\footnotesize,
  every node/.style={align=center},
  ag/.style={draw=SkyBlue!65!black, fill=SkyBlue!13, rounded corners=2.5pt,
             inner sep=4pt, font=\scriptsize, minimum height=16mm,
             line width=0.5pt, anchor=north},
  adv/.style={draw=BurntOrange!72!black, fill=BurntOrange!12, rounded corners=2.5pt,
              inner sep=5pt, font=\scriptsize, line width=0.6pt},
  lead/.style={draw=black!60, fill=black!5, rounded corners=2.5pt, inner sep=5pt,
               font=\scriptsize, line width=0.7pt},
  det/.style={draw=SeaGreen!60!black, fill=SeaGreen!11, rounded corners=2.5pt,
              inner sep=5pt, font=\scriptsize, line width=0.55pt},
  own/.style={draw=RoyalPurple!60!black, fill=RoyalPurple!10, rounded corners=2.5pt,
              inner sep=5pt, font=\scriptsize, line width=0.55pt},
  fl/.style={-{Latex[length=2mm,width=1.5mm]}, draw=black!55, line width=0.6pt},
  atk/.style={-{Latex[length=2mm,width=1.5mm]}, draw=BurntOrange!72!black,
              line width=0.6pt, dashed},
  note/.style={font=\scriptsize\itshape, text=black!60},
  hd/.style={font=\scriptsize\bfseries, text=black!55},
]

% ---------- the seven domain agents ----------
\node[hd, anchor=west] at (-7.5,3.5)
  {\textsc{propose} --- seven domain agents, 43 gaps, one owner each};
\foreach \i/\name/\n in {
  0/{\code{AGENT\_1}\\formula, cross-Veda}/3,
  1/{\code{AGENT\_2}\\entity, deity identity}/12,
  2/{\code{AGENT\_3}\\morphology, semantics}/7,
  3/{\code{AGENT\_4}\\ritual}/4,
  4/{\code{AGENT\_5}\\quality, reference}/6,
  5/{\code{AGENT\_6}\\product surface}/5,
  6/{\code{AGENT\_7}\\translation, S\=amaveda}/6} {
  \node[ag, text width=18mm] (a\i) at ({-7.0+\i*2.35},3.05)
    {\name\\[-1pt]\tiny \n\ gaps};
}
\node[note, anchor=west, text width=152mm, fill=white, inner sep=2pt] at (-7.5,0.16)
  {Each emits measurements, row-level evidence and a
   \code{proposed\_mutations.json}. \textbf{None writes to the graph.}};

% ---------- the adversary ----------
\node[adv, text width=52mm, anchor=north] (ag8) at (-4.15,-0.78)
  {\textbf{\code{AGENT\_8}} --- no gaps assigned\\[1pt]
   \tiny database access recorded \code{READ\_ONLY}\\[2pt]
   \tiny\raggedright
   Found, among others: a registry check that ``never read the status field'';
   an export under-declaring 31 of 36 labels, hiding 129{,}691 edges;
   5{,}123 search rows carrying a private-use sentinel, so a reader's query
   returned 0 where the sentinel form returned 94.};

\node[hd, anchor=west] at (-7.5,-0.46) {\textsc{attack}};

\foreach \i in {0,...,6} {
  \draw[atk] (a\i.south) -- ++(0,-0.95) -| ($(ag8.north)+({-1.2+\i*0.4},0)$);
}

% ---------- the lead ----------
\node[lead, text width=52mm, anchor=north] (lead) at (3.05,-0.78)
  {\textbf{\code{lead}} --- integrator\\[2pt]
   \tiny\raggedright
   Re-measures before ruling: each ruling records a
   \code{lead\_measurement} separately from the agent's claim. Four rulings
   survive; one corrects an agent's own figure (4{,}865 $\to$ 35{,}139 after its
   own staged delta), and one records a \emph{trade} --- closing one gap deepens
   two others --- rather than reporting the improvement alone.};
\node[hd, anchor=west] at (0.55,-0.46) {\textsc{adjudicate}};

\draw[fl] (ag8.east) -- node[note, above, font=\tiny, pos=0.5, align=center]
  {conflicts,\\measured} (lead.west);
\draw[fl] (a6.south east) -- ++(0.55,0) |- node[note, above, font=\tiny, pos=0.72]
  {proposals} ($(lead.north)+(1.1,0)$);

% ---------- worked example ----------
\node[draw=black!25, rounded corners=3pt, dashed, inner sep=6pt, line width=0.5pt,
      fill=black!2, text width=152mm, font=\scriptsize, align=left]
  (ex) at (-0.55,-5.15)
  {\textbf{One ruling, end to end (\code{R2\_ASSERTION\_ROLE\_REANCHOR}).}\quad
   \textsc{propose}: Agent~3 asks to re-anchor \code{ASSERTION\_ROLE} so role
   fillers hang from semantic assertions.\quad
   \textsc{attack}: Agent~8 measures that this creates \textbf{2{,}052 endpoint
   signature violations} --- the declared signature expects a passage endpoint and
   \emph{0} of the assertion nodes carry that label.\quad
   \textsc{verify}: the lead reproduces the measurement independently.\quad
   \textsc{rule}: not approved and not rejected, but \emph{approved with the
   ontology moving in the same change} --- the migration card had declared the new
   signature all along, and either half alone leaves 2{,}052 violations. A test now
   pins the signature to the card.};

% ---------- the two things agents cannot do ----------
\node[det, text width=64mm] (det) at (-4.0,-8.05)
  {\textbf{deterministic code decides promotion}\\[2pt]
   \tiny\raggedright a frozen set of trust classes may be written
   \code{ACCEPTED} without review; the model's class is not in it, and the
   constructor raises. Acceptance additionally requires a predicate to have been
   \emph{unlocked} by a measured precision --- the unlocked set is empty by
   construction.};
\node[own, text width=64mm] (own) at (3.4,-8.05)
  {\textbf{the owner decides scope and \textsc{go}}\\[2pt]
   \tiny\raggedright every mutation runs plan $\to$ dry run $\to$ owner
   \textsc{go} $\to$ execute $\to$ readback against a hashed plan. Owner rulings
   are cited by identifier from the closure tests they license, so a test cannot
   silently redefine what it checks.};

\draw[fl] (ex.south) -- ++(0,-0.35) -| (det.north);
\draw[fl] (ex.south) -- ++(0,-0.35) -| (own.north);

\node[note, text width=152mm, align=left] at (-0.55,-9.55)
  {The agents are not statistically independent: they share a model family, so a
   correlated blind spot would be invisible to the whole arrangement. What the
   artefacts show is only that a \emph{differently instructed} pass finds what an
   earlier pass missed --- repeatedly, and including in its own prior work.};
\end{tikzpicture}
}
\caption{Propose $\to$ attack $\to$ adjudicate, using the repository's own agent
names and counts. Seven domain agents own 43 gaps between them and write no graph
mutations; an eighth agent has no gaps and read-only database access, and audits the
other seven; a lead re-measures before ruling. The boxed example is one ruling
end to end. The caveat at the foot is not decorative: these agents share a model
family and are not statistically independent.}
\label{fig:adversarial}
\end{figure}

\subsection{Five hostile rounds}

Before release, five further adversarial rounds ran over the candidate
(\cref{sec:failure}). The pattern to report is that later rounds found defects in
earlier rounds' work, including one round that closed two critical findings its own
hostile pass had raised against itself. The first round's finding was structural:
30{,}266 semantic-assertion nodes carried no assertion identifier and fell through to
their source passage's canonical key, collapsing two public objects onto one and
silently dropping assertion-to-passage edges as self-edges.

We are careful about what this shows. It shows that a differently-instructed pass
finds what an earlier pass missed. It does not show statistical independence: the
passes share a model family, and a correlated blind spot would be invisible to all
of them (\cref{sec:threats}).

\subsection{Three gates per domain}
\label{sec:gates}

Fourteen domains passed through a three-gate structure, adopted after four of four
domains failed a third gate that had not previously existed. Gate~A verifies the
artefact's own integrity --- that its manifest checksums hold. Gate~B verifies the
rows against the live graph: coverage, addressing, provenance, and the central
importability policy. Gate~C attacks. In the case we treat at length
(\cref{sec:case-samaveda}) Gate~C was implemented in a separate file from Gate~B
explicitly ``so the same codepath cannot validate itself'', named the fields it
refused to read, re-implemented the normalisation from prose rather than reusing the
pipeline's code, and carried a negative control.

The owner-decision chain that produced Gate~C is itself part of the method. A
condition required that ``existing validation gates pass''; the eligibility record
showed one gate unknown and the other not run. The condition was recorded as
\emph{unmet}, with the reasoning that reading ``gates'' as the one gate that already
existed would clear the board by narrowing a word, and that the conservative reading
is the one that cannot manufacture a closure. The owner confirmed the conservative
reading, the gates were built, and the new gate landed a real defect on its first
run.

\subsection{Model output is confined by mechanism}
\label{sec:mechanism}

\Cref{sec:model-policy} states the policy; here we note where it bites. The
extractor's validator rejects with one code per failure mode, and its central check
is whether a cited token was present in the packet the model was shown --- a
fabricated citation is decidable and is discarded. The predicate vocabulary is
closed, and five predicates a model would naturally reach for are members
\emph{so that proposing one is refused by name with a recorded reason}. Acceptance
requires a predicate to have been unlocked by a measured precision against gold; the
unlocked set is empty by construction and remains empty, so no assertion has ever
been auto-accepted.

Runtime provenance is recorded rather than assumed, on a rule adopted after a
concrete failure: earlier pilot output had been recorded as model-generated when it
was in fact produced by regular expressions. The decision record calls that
``provenance that looked like evidence and was not'', and notes that a wrong runtime
string is the same defect in one field. Self-reported model identity is classified as
unattested, and provider-attested build metadata is marked unavailable because no
provider supplies it.

\subsection{Who did what}

% Hand-authored: this table records a division of responsibility read off the
% artefacts, which is a judgement about what the artefacts show, not a measurement.
% Each row is supported by the evidence cited in Appendix C.
\begin{table}[tb]
\centering
\caption{Division of responsibility by construction stage.
\textbullet\ = performs the work; \textopenbullet\ = participates but cannot decide;
--- = not involved. The rows that determine what the graph \emph{asserts} ---
identity, promotion, mutation and certification --- have no filled agent cell.}
\label{tab:responsibility}
\small
\begin{tabularx}{\linewidth}{@{}>{\raggedright\arraybackslash}p{40mm}ccc
                             >{\raggedright\arraybackslash}X@{}}
\toprule
Stage & \rotatebox{0}{Det.\ code} & \rotatebox{0}{Agent} & \rotatebox{0}{Owner}
      & Where the decision is recorded \\
\midrule

Source discovery and rights & \textopenbullet & \textbullet & \textbullet
  & rights registry, per file; owner rules on conflicting statements \\

Snapshot and hashing & \textbullet & --- & ---
  & content-addressed store, 11 required sidecar fields \\

Canonical parse and identity & \textbullet & --- & ---
  & \mbox{UUIDv5} over a \mbox{URN}; no model in the path \\

Ontology and vocabulary design & \textopenbullet & \textbullet & \textbullet
  & decision records; closed enums with named refusals \\

Candidate generation & \textbullet & \textbullet & ---
  & staging rows; deterministic rules and model proposals both land here \\

Validation of candidates & \textbullet & --- & ---
  & 20 rejection codes, one code per failure mode \\

Adversarial audit & \textopenbullet & \textbullet & ---
  & a read-only agent with no gaps assigned; findings are counts, re-runnable \\

Reconciling conflicting proposals & \textopenbullet & \textbullet & \textbullet
  & lead rulings, each recording its own independent measurement \\

Promotion candidate $\to$ accepted & \textbullet & --- & \textopenbullet
  & a frozen set of trust classes; the model's class is not in it \\

Graph mutation & \textbullet & --- & \textbullet
  & plan $\to$ dry run $\to$ owner \textsc{go} $\to$ execute $\to$ readback \\

Readback and invariants & \textbullet & --- & ---
  & rows sent vs rows landed; 20 invariants, 11 scorecard gates \\

Scope and withholding rulings & --- & \textopenbullet & \textbullet
  & owner decisions, cited by identifier from the closure tests they license \\

Release certification & \textbullet & \textopenbullet & \textbullet
  & certification document; withheld once on a single blocker \\

Manuscript drafting & --- & \textbullet & \textbullet
  & this paper; claims and interpretation are the author's \\

\bottomrule
\end{tabularx}
\end{table}


The distribution in \cref{tab:responsibility} is the answer to RQ2. Agents are
present in almost every row and decisive in none of the rows that define truth:
identity, promotion, graph mutation and release certification are deterministic code
or a human decision, every time.
